\documentclass[12pt,a4paper]{article}

% ============================================
% PACKAGES
% ============================================
\usepackage[utf8]{inputenc}
\usepackage[T1]{fontenc}
\usepackage[french]{babel}
\usepackage{amsmath, amsfonts}
\usepackage{graphicx}
\usepackage{float}
\usepackage{booktabs}
\usepackage{geometry}
\usepackage{xcolor}
\usepackage{listings}
\usepackage{hyperref}
\usepackage{fancyhdr}

% ============================================
% CONFIGURATION
% ============================================
\geometry{left=2.5cm, right=2.5cm, top=2.5cm, bottom=2.5cm}
\hypersetup{colorlinks=true, linkcolor=blue, urlcolor=cyan}

% Configuration des listings pour R
\lstdefinestyle{Rstyle}{
    language=R,
    basicstyle=\ttfamily\footnotesize,
    keywordstyle=\color{blue}\bfseries,
    commentstyle=\color{green!60!black},
    numbers=left,
    numberstyle=\tiny\color{gray},
    backgroundcolor=\color{gray!10},
    frame=single,
    breaklines=true,
    tabsize=2,
    literate={é}{{\'e}}1 {è}{{\`e}}1 {ê}{{\^e}}1 {à}{{\`a}}1 {ù}{{\`u}}1
}

% Configuration des listings pour Python
\lstdefinestyle{Pythonstyle}{
    language=Python,
    basicstyle=\ttfamily\footnotesize,
    keywordstyle=\color{blue}\bfseries,
    commentstyle=\color{green!60!black},
    numbers=left,
    numberstyle=\tiny\color{gray},
    backgroundcolor=\color{yellow!10},
    frame=single,
    breaklines=true,
    tabsize=4,
    literate={é}{{\'e}}1 {è}{{\`e}}1 {ê}{{\^e}}1 {à}{{\`a}}1 {ù}{{\`u}}1
}

% En-tête
\pagestyle{fancy}
\fancyhf{}
\rhead{GMM2 - Exercice 4}
\lhead{Analyse Factorielle des Données Mixtes}
\cfoot{\thepage}

% ============================================
% DOCUMENT
% ============================================
\begin{document}

% ============================================
% TITRE
% ============================================
\begin{center}
    \rule{\linewidth}{0.5mm}\\[0.4cm]
    {\Huge \textbf{EXERCICE 4}}\\[0.3cm]
    {\Large \textbf{Analyse Factorielle des Données Mixtes (AFDM)}}\\[0.2cm]
    {\large Base de données : Bank}\\[0.1cm]
    \rule{\linewidth}{0.5mm}\\[1cm]
\end{center}

% ============================================
% PRÉSENTATION DES DONNÉES
% ============================================
\section{Présentation des Données}

\subsection{Description}

La base de données \textbf{Bank} contient des informations sur des clients d'une banque. Elle est caractérisée par la présence de \textbf{variables mixtes} : à la fois quantitatives et qualitatives, ce qui justifie l'utilisation de l'AFDM.

\begin{itemize}
    \item \textbf{Nombre d'observations} : 8 clients
    \item \textbf{Nombre de variables} : 7 variables
    \item \textbf{Variables quantitatives} : 2 (age, balance)
    \item \textbf{Variables qualitatives} : 5 (job, marital, education, housing, loan)
\end{itemize}

\subsection{Description des variables}

\begin{table}[H]
\centering
\caption{Description des variables de la base Bank}
\begin{tabular}{|l|c|l|}
\hline
\textbf{Variable} & \textbf{Type} & \textbf{Description / Modalités} \\
\hline
age & Quantitative & Âge du client (en années) \\
balance & Quantitative & Solde du compte bancaire (en euros) \\
job & Qualitative & Profession (admin, technician, management, services, student, retired) \\
marital & Qualitative & Statut marital (single, married, divorced) \\
education & Qualitative & Niveau d'éducation (secondary, tertiary) \\
housing & Qualitative & Crédit immobilier (yes, no) \\
loan & Qualitative & Prêt personnel (yes, no) \\
\hline
\end{tabular}
\end{table}

\subsection{Tableau des données}

\begin{table}[H]
\centering
\caption{Données brutes de la base Bank}
\footnotesize
\begin{tabular}{|c|c|l|l|l|c|c|c|}
\hline
\textbf{Ind.} & \textbf{Age} & \textbf{Job} & \textbf{Marital} & \textbf{Education} & \textbf{Balance} & \textbf{Housing} & \textbf{Loan} \\
\hline
1 & 30 & admin & single & secondary & 1200 & yes & no \\
2 & 45 & technician & married & tertiary & 3500 & yes & no \\
3 & 50 & management & married & tertiary & 8000 & no & no \\
4 & 35 & services & single & secondary & 1500 & yes & yes \\
5 & 28 & student & single & tertiary & 500 & no & no \\
6 & 60 & retired & married & secondary & 4000 & no & no \\
7 & 40 & admin & divorced & secondary & 2300 & yes & yes \\
8 & 33 & services & married & secondary & 1800 & yes & no \\
\hline
\end{tabular}
\end{table}

\subsection{Statistiques descriptives}

\subsubsection{Variables quantitatives}

\begin{table}[H]
\centering
\caption{Statistiques descriptives des variables quantitatives}
\begin{tabular}{|l|c|c|c|c|c|}
\hline
\textbf{Variable} & \textbf{Moyenne} & \textbf{Écart-type} & \textbf{Min} & \textbf{Médiane} & \textbf{Max} \\
\hline
age & 40.12 & 10.97 & 28 & 37.5 & 60 \\
balance & 2850 & 2382.68 & 500 & 2050 & 8000 \\
\hline
\end{tabular}
\end{table}

\textbf{Observations :}
\begin{itemize}
    \item L'âge moyen des clients est de 40 ans, avec une dispersion importante (écart-type de 11 ans)
    \item Le solde moyen est de 2850€, avec une forte variabilité (de 500€ à 8000€)
\end{itemize}

\subsubsection{Variables qualitatives}

\begin{table}[H]
\centering
\caption{Répartition des modalités des variables qualitatives}
\begin{tabular}{|l|l|c|c|}
\hline
\textbf{Variable} & \textbf{Modalité} & \textbf{Effectif} & \textbf{Pourcentage} \\
\hline
\multirow{6}{*}{job} & admin & 2 & 25\% \\
 & services & 2 & 25\% \\
 & technician & 1 & 12.5\% \\
 & management & 1 & 12.5\% \\
 & student & 1 & 12.5\% \\
 & retired & 1 & 12.5\% \\
\hline
\multirow{3}{*}{marital} & married & 4 & 50\% \\
 & single & 3 & 37.5\% \\
 & divorced & 1 & 12.5\% \\
\hline
\multirow{2}{*}{education} & secondary & 5 & 62.5\% \\
 & tertiary & 3 & 37.5\% \\
\hline
\multirow{2}{*}{housing} & yes & 5 & 62.5\% \\
 & no & 3 & 37.5\% \\
\hline
\multirow{2}{*}{loan} & no & 6 & 75\% \\
 & yes & 2 & 25\% \\
\hline
\end{tabular}
\end{table}

\textbf{Observations :}
\begin{itemize}
    \item \textbf{Job} : Diversité des professions, avec prédominance des administratifs et des employés de services
    \item \textbf{Marital} : Moitié des clients sont mariés
    \item \textbf{Education} : Majorité de niveau secondaire (62.5\%)
    \item \textbf{Housing} : La majorité des clients (62.5\%) ont un crédit immobilier
    \item \textbf{Loan} : Seuls 25\% des clients ont un prêt personnel
\end{itemize}

% ============================================
% MÉTHODOLOGIE
% ============================================
\section{Méthodologie}

L'\textbf{Analyse Factorielle des Données Mixtes (AFDM)} est une méthode d'analyse multivariée adaptée aux jeux de données contenant simultanément des \textbf{variables quantitatives} et \textbf{qualitatives}. Elle combine les principes de :

\begin{itemize}
    \item L'\textbf{ACP} (Analyse en Composantes Principales) pour les variables quantitatives
    \item L'\textbf{ACM} (Analyse des Correspondances Multiples) pour les variables qualitatives
\end{itemize}

\textbf{Principe :} L'AFDM équilibre la contribution des deux types de variables en :
\begin{enumerate}
    \item \textbf{Normalisant} les variables quantitatives (centrage-réduction)
    \item \textbf{Pondérant} les variables qualitatives selon leur nombre de modalités
    \item \textbf{Maximisant} l'inertie expliquée sur des axes factoriels communs
\end{enumerate}

\textbf{Objectifs :}
\begin{enumerate}
    \item Visualiser les relations entre individus
    \item Identifier les associations entre variables de types différents
    \item Détecter des profils types de clients
    \item Réduire la dimensionnalité tout en conservant l'information mixte
\end{enumerate}

% ============================================
% RÉSULTATS
% ============================================
\section{Résultats}

\subsection{Valeurs propres et inertie}

\begin{table}[H]
\centering
\caption{Valeurs propres et variance expliquée}
\begin{tabular}{|c|c|c|c|}
\hline
\textbf{Dimension} & \textbf{Valeur propre} & \textbf{Variance (\%)} & \textbf{Cumul (\%)} \\
\hline
Dim.1 & 8.694 & 35.81 & 35.81 \\
Dim.2 & 5.129 & 21.12 & 56.93 \\
Dim.3 & 3.915 & 16.12 & 73.05 \\
Dim.4 & 3.737 & 15.39 & 88.45 \\
Dim.5 & 2.806 & 11.55 & 100.00 \\
\hline
\end{tabular}
\end{table}

\textbf{Observations :}
\begin{itemize}
    \item Les \textbf{deux premières dimensions} expliquent \textbf{56.93\%} de l'inertie totale
    \item Les \textbf{trois premières dimensions} atteignent \textbf{73.05\%}
    \item La première dimension est la plus importante (35.81\%), ce qui suggère une structure principale forte
\end{itemize}

\subsection{Résultats numériques}

\begin{figure}[H]
    \centering
    \includegraphics[width=0.95\textwidth]{"resultats numeriques.png"}
    \caption{Résultats numériques de l'AFDM (valeurs propres, coordonnées, contributions)}
    \label{fig:resultats}
\end{figure}

\subsection{Graphiques de l'AFDM}

\begin{figure}[H]
    \centering
    \includegraphics[width=0.9\textwidth]{"FADM_plot variables.png"}
    \caption{Projection des variables sur le plan factoriel (Dim 1 - Dim 2)}
    \label{fig:variables}
\end{figure}

\subsection{Contributions des variables}

\begin{figure}[H]
    \centering
    \includegraphics[width=0.85\textwidth]{"contibution dim1 variables.png"}
    \caption{Contribution des variables à la dimension 1}
    \label{fig:contrib_var1}
\end{figure}

\begin{figure}[H]
    \centering
    \includegraphics[width=0.85\textwidth]{"contibution dim2 variables.png"}
    \caption{Contribution des variables à la dimension 2}
    \label{fig:contrib_var2}
\end{figure}

\subsection{Contributions des individus}

\begin{figure}[H]
    \centering
    \includegraphics[width=0.85\textwidth]{"contibution dim1 individus.png"}
    \caption{Contribution des individus à la dimension 1}
    \label{fig:contrib_ind1}
\end{figure}

\begin{figure}[H]
    \centering
    \includegraphics[width=0.85\textwidth]{"contibution dim2 individus.png"}
    \caption{Contribution des individus à la dimension 2}
    \label{fig:contrib_ind2}
\end{figure}

\subsection{Contributions numériques des variables}

\begin{table}[H]
\centering
\caption{Contributions des variables aux deux premiers axes (\%)}
\begin{tabular}{|l|c|c|}
\hline
\textbf{Variable} & \textbf{Dim.1 (\%)} & \textbf{Dim.2 (\%)} \\
\hline
age & 3.45 & 10.66 \\
balance & 5.30 & 5.71 \\
education & 5.92 & 3.29 \\
housing & 6.79 & 0.57 \\
job & \textbf{10.68} & \textbf{17.18} \\
loan & 6.35 & 0.50 \\
marital & 6.77 & 11.91 \\
\hline
\end{tabular}
\end{table}

% ============================================
% INTERPRÉTATION
% ============================================
\section{Interprétation des Résultats}

\subsection{Interprétation des axes factoriels}

\subsubsection{Axe 1 : Profil socio-économique global (35.81\% d'inertie)}

Le premier axe factoriel représente un \textbf{gradient socio-économique} qui oppose :

\textbf{Côté positif :}
\begin{itemize}
    \item \textbf{Clients plus âgés} avec des \textbf{soldes bancaires élevés}
    \item Professions de type \textbf{management} et \textbf{retired}
    \item Statut \textbf{married}
    \item Niveau d'éducation \textbf{tertiary}
    \item Pas de crédit immobilier ni de prêt personnel
\end{itemize}

\textbf{Côté négatif :}
\begin{itemize}
    \item \textbf{Clients plus jeunes} avec des \textbf{soldes modestes}
    \item Professions de type \textbf{services}, \textbf{admin} ou \textbf{student}
    \item Statut \textbf{single}
    \item Niveau \textbf{secondary}
    \item Présence de crédits (housing = yes)
\end{itemize}

\textbf{Interprétation :} Cet axe distingue les \textbf{clients aisés et établis} (âgés, bonne situation financière, mariés, niveau d'éducation supérieur) des \textbf{clients plus modestes et en début de vie active} (jeunes, soldes faibles, célibataires, niveau secondaire, endettés).

La variable \textbf{job} contribue le plus à cet axe (10.68\%), suivie de \textbf{marital} (6.77\%) et \textbf{housing} (6.79\%).

\subsubsection{Axe 2 : Profession et situation familiale (21.12\% d'inertie)}

Le deuxième axe apporte une information complémentaire en distinguant :

\textbf{Côté positif :}
\begin{itemize}
    \item Profession \textbf{student}
    \item Clients \textbf{jeunes} 
    \item Statut \textbf{single}
\end{itemize}

\textbf{Côté négatif :}
\begin{itemize}
    \item Professions \textbf{retired} et \textbf{management}
    \item Clients plus \textbf{âgés}
    \item Statut \textbf{married}
\end{itemize}

\textbf{Interprétation :} L'axe 2 oppose les \textbf{étudiants} (jeunes, célibataires) aux \textbf{retraités et cadres} (âgés, mariés). Il capture la dimension du \textbf{cycle de vie}.

Les principales contributions viennent de \textbf{job} (17.18\%), \textbf{marital} (11.91\%) et \textbf{age} (10.66\%).

\subsection{Analyse des profils types}

\subsubsection{Profil 1 : Clients aisés et établis}

Ce profil regroupe les individus caractérisés par :
\begin{itemize}
    \item Âge élevé (50-60 ans)
    \item Profession de niveau supérieur (management, retired)
    \item Marié(e)
    \item Niveau d'éducation tertiary
    \item Solde bancaire élevé (4000-8000€)
    \item Pas de prêts en cours
\end{itemize}

\textbf{Exemple} : Individu 3 (management, 50 ans, married, tertiary, balance=8000€)

\subsubsection{Profil 2 : Jeunes actifs avec crédits}

Ce profil est caractérisé par :
\begin{itemize}
    \item Âge moyen (30-40 ans)
    \item Profession intermédiaire (admin, services, technician)
    \item Célibataire ou marié(e)
    \item Niveau secondary
    \item Solde modeste (1200-2300€)
    \item Crédit immobilier et/ou prêt personnel
\end{itemize}

\textbf{Exemples} : Individus 1, 4, 7 (admin/services, 30-40 ans, secondary, housing=yes)

\subsubsection{Profil 3 : Étudiants}

Ce profil particulier se caractérise par :
\begin{itemize}
    \item Très jeune (28 ans)
    \item Profession : student
    \item Célibataire
    \item Niveau tertiary (études supérieures en cours)
    \item Solde très faible (500€)
    \item Pas de crédits
\end{itemize}

\textbf{Exemple} : Individu 5 (student, 28 ans, single, tertiary, balance=500€)

\subsection{Relations entre variables quantitatives et qualitatives}

L'AFDM révèle des associations intéressantes :

\begin{enumerate}
    \item \textbf{Age et marital} : Les clients âgés sont majoritairement mariés, tandis que les jeunes sont célibataires. Cette corrélation logique reflète le cycle de vie familial.
    
    \item \textbf{Balance et job} : Les professions de niveau supérieur (management, retired) sont associées à des soldes élevés, confirmant le lien entre statut professionnel et capacité d'épargne.
    
    \item \textbf{Education et job} : Le niveau tertiary est associé aux professions qualifiées (management, technician), tandis que secondary est lié aux emplois administratifs et de services.
    
    \item \textbf{Housing et age} : Les clients jeunes et intermédiaires sont plus susceptibles d'avoir un crédit immobilier, correspondant à la phase d'acquisition du logement.
    
    \item \textbf{Loan et housing} : Les prêts personnels sont souvent combinés avec un crédit immobilier, suggérant un profil d'endettement multiple.
\end{enumerate}

\subsection{Synthèse de l'interprétation}

L'AFDM appliquée aux données Bank révèle :

\begin{enumerate}
    \item \textbf{Une structure socio-économique claire} (Axe 1) : Opposition entre clients aisés/établis et clients modestes/en début de carrière.
    
    \item \textbf{Un effet du cycle de vie} (Axe 2) : Distinction entre les phases de vie (étudiant, actif, retraité) avec leurs caractéristiques associées.
    
    \item \textbf{Des profils clients cohérents} : L'analyse identifie au moins trois segments distincts exploitables pour du ciblage marketing ou la gestion du risque crédit.
    
    \item \textbf{L'importance de la profession} : La variable job contribue le plus aux deux axes, confirmant son rôle central dans la caractérisation des clients.
    
    \item \textbf{La complémentarité des variables} : Les variables quantitatives (age, balance) et qualitatives (job, marital, education) apportent des informations convergentes sur les profils clients.
\end{enumerate}

% ============================================
% CODES
% ============================================
\section{Codes R et Python}

\subsection{Code R}

\begin{lstlisting}[style=Rstyle, caption={Code R pour l'AFDM avec FactoMineR et factoextra}]
# Chargement des librairies necessaires
library(readxl)
library(FactoMineR)
library(ggplot2)
library(factoextra)
library(PCAmixdata)

# Importation des donnees
bank <- read_excel("bank.xlsx")

# Exploration des donnees
dim(bank)          # Dimensions du jeu de donnees
head(bank, 10)     # Apercu des 10 premieres lignes
str(bank)          # Structure des variables

# Statistiques descriptives
summary(bank)

# Realisation de l'AFDM
res.famd <- FAMD(bank, graph = FALSE)

# Affichage des resultats numeriques
summary(res.famd)

# Graphique des variables
fviz_famd_var(res.famd)

# Eboulis des valeurs propres
fviz_eig(res.famd)

# Contribution des variables a la dimension 1
fviz_contrib(res.famd, choice = "var", axes = 1)

# Contribution des variables a la dimension 2
fviz_contrib(res.famd, choice = "var", axes = 2)

# Contribution des individus a la dimension 1
fviz_contrib(res.famd, choice = "ind", axes = 1)

# Contribution des individus a la dimension 2
fviz_contrib(res.famd, choice = "ind", axes = 2)
\end{lstlisting}

\subsection{Code Python}

\begin{lstlisting}[style=Pythonstyle, caption={Code Python pour l'AFDM avec prince}]
# Chargement des librairies
import pandas as pd
import numpy as np
import matplotlib.pyplot as plt
import prince

# Chargement des donnees
bank = pd.read_excel("bank.xlsx")

# Conversion des types pour prince
# Variables quantitatives en float
bank['age'] = bank['age'].astype(float)
bank['balance'] = bank['balance'].astype(float)

# Variables qualitatives en category
for col in ['job', 'marital', 'education', 'housing', 'loan']:
    bank[col] = bank[col].astype('category')

# Affichage des informations
print("Dimensions :", bank.shape)
print("\nApercu :")
print(bank.head())

# Statistiques descriptives
print("\nStatistiques quantitatives :")
print(bank.describe())

# Realisation de l'AFDM
famd = prince.FAMD(
    n_components=5,
    n_iter=3,
    random_state=42
)
famd = famd.fit(bank)

# Valeurs propres
eigenvalues = famd.eigenvalues_
total_inertia = sum(eigenvalues)
print("\n=== VALEURS PROPRES ===")
for i, eig in enumerate(eigenvalues):
    var_pct = eig / total_inertia * 100
    cum_pct = sum(eigenvalues[:i+1]) / total_inertia * 100
    print(f"Dim.{i+1}: {eig:.3f} ({var_pct:.2f}%, cumul: {cum_pct:.2f}%)")

# Coordonnees des individus
print("\n=== COORDONNEES DES INDIVIDUS ===")
row_coords = famd.row_coordinates(bank)
print(row_coords.round(3))

# Coordonnees des variables
print("\n=== COORDONNEES DES VARIABLES ===")
col_coords = famd.column_coordinates(bank)
print(col_coords.round(3))

# Contributions des variables
print("\n=== CONTRIBUTIONS DES VARIABLES (%) ===")
var_contrib = famd.column_contributions_ * 100
print(var_contrib.round(2))

# Visualisation du biplot
ax = famd.plot(
    bank,
    x_component=0,
    y_component=1,
    show_column_markers=True,
    show_row_markers=True
)
plt.title("AFDM - Biplot Dimensions 1 et 2")
plt.tight_layout()
plt.savefig("famd_biplot.png", dpi=150)
plt.show()
\end{lstlisting}

% ============================================
% CONCLUSION
% ============================================
\section{Conclusion}

L'Analyse Factorielle des Données Mixtes appliquée aux données Bank a permis de :

\begin{enumerate}
    \item \textbf{Traiter simultanément les deux types de variables} : Les 2 variables quantitatives (age, balance) et les 5 variables qualitatives (job, marital, education, housing, loan) ont été intégrées dans une analyse unique.
    
    \item \textbf{Identifier deux dimensions structurantes} : L'axe 1 (35.81\%) capture le niveau socio-économique global, tandis que l'axe 2 (21.12\%) reflète le cycle de vie des clients.
    
    \item \textbf{Révéler des profils clients distincts} : Trois segments ont été identifiés (clients aisés, jeunes actifs endettés, étudiants), exploitables pour des stratégies commerciales différenciées.
    
    \item \textbf{Mettre en évidence l'importance de la profession} : La variable job s'avère la plus discriminante, contribuant fortement aux deux axes principaux.
    
    \item \textbf{Confirmer des associations logiques} : Les relations observées (âge-statut marital, profession-revenus, éducation-emploi) sont cohérentes avec les connaissances sociologiques.
\end{enumerate}

\textbf{Limites :} La petite taille de l'échantillon (8 observations) limite la généralisation des résultats. Une analyse sur un échantillon plus représentatif permettrait de valider ces profils.

\textbf{Applications potentielles :}
\begin{itemize}
    \item Segmentation marketing de la clientèle bancaire
    \item Scoring de risque crédit
    \item Personnalisation des offres de produits financiers
\end{itemize}

\end{document}
